% Canonical LaTeX chapter; migrated from docs/api.md.
\hypertarget{http-api}{%
\subsection{HTTP API}\label{http-api}}

The server publishes its precise OpenAPI schema at
\texttt{/openapi.json} and interactive reference at \texttt{/docs}. JSON
business-validation failures return \textbf{400} with a \texttt{detail}
string; malformed request schemas return \textbf{422}; missing
authentication returns \textbf{401}; missing pack resources return
\textbf{404}; write conflicts return \textbf{409}. Private API responses
have \texttt{Cache-Control: no-store}.

In \texttt{AUTH\_MODE=multi}, use
\texttt{Authorization: Bearer \textless{}token\textgreater{}}. The
health endpoint and login are unauthenticated. Accounts are provisioned
using \texttt{ascend-server user-add}. There is no public registration
endpoint. Passwords use scrypt; tokens are hashed at rest. Browser
tokens are session-scoped, and CLI tokens are stored with mode 0600.

\begin{longtable}[]{@{}P{.45\textwidth}P{.50\textwidth}@{}}
\toprule\noalign{}
Route & Contract \\
\midrule\noalign{}
\endhead
\bottomrule\noalign{}
\endlastfoot
\path|GET /api/health| & \texttt{\{status, auth\_mode\}} \\
\path|POST /api/login| & \texttt{\{username,password\}} →
\texttt{\{token\}} \\
\path|GET /api/me?minutes=30| & Mission, next/quick-win node data,
current level, streak and today state \\
\path|GET /api/levels?pack=<id>| & Levels with
node states, active days and minutes \\
\path|GET /api/challenges/{reference}| & Caller-scoped node, level, session,
public exercise and optional exercise error. Read-only recovery; does not start timing \\
\path|POST /api/entries| &
Required: \texttt{learned}, \texttt{next\_step}. Optional: \texttt{node\_id}, \texttt{proof}, \texttt{minutes}, \texttt{complete}, \texttt{force}, \texttt{mission\_id}.
→ entry ID, logical day and completion result \\
\path|GET /api/entries| & Caller-scoped learning entries \\
\path|POST /api/nodes/{pack:node}/mark| &
\texttt{\{status,note?\}}; status is done/skipped/todo \\
\path|POST /api/packs| & Multipart \texttt{file}: UTF-8 TOML, maximum
2 MB; lint and install \\
\path|GET /api/packs| & Installed pack IDs and review/trust
metadata \\
\path|GET /api/packs/{id}| & Pack metadata and preview nodes \\
\path|GET /api/stats| & Derived streaks, days, freeze allowance and
eight weekly minute totals \\
\path|GET /api/heatmap| & Day/count values for the last 12 calendar
weeks \\
\path|GET /api/search?q=text| & Reverse chronological substring
matches \\
\path|GET /api/account| & Current account and reminder preferences \\
\path|POST /api/command| &
\texttt{\{argv:{[}...{]},pack\_text?:"..."\}}; same application commands
as the CLI \\
\path|POST /api/logout| & Revoke the current token \\
\path|GET /api/push/key| & Configured public VAPID key \\
\path|POST /api/push/subscriptions| & Browser PushSubscription JSON;
user-scoped upsert \\
\path|DELETE /api/push/subscriptions| & \texttt{\{endpoint\}}; remove
caller\textquotesingle s subscription \\
\path|POST /api/telegram/webhook| & Telegram Update; requires the
Telegram webhook secret header \\
\path|GET /api/public/{username}/streak| & Minimal opt-in streak
data; 404 otherwise \\
\path|GET /api/public/{username}/badge.svg| & Opt-in SVG badge \\
\end{longtable}

Overview, level and pack listings return preview nodes. They omit full task
instructions, materials, hints, completion criteria and verifier internals.
Resource links, labels and notes remain available before Start.
The challenge endpoint sets \texttt{details\_available} and returns
full guidance only for an existing caller-owned session or historical completion.
An ended attempt retains review access. \texttt{has\_attempt} lets the interface
offer Review attempt without starting another timer. The \texttt{today} and
\texttt{levels} command responses follow the same preview contract.

\texttt{/api/command} supports mission lifecycle, config, edit/undo,
backfill, split, verification, and parking. It rejects client-only timer
installation and pack scaffolding. Pack paths sent by a client must
include \texttt{pack\_text}; the server does not interpret them as
arbitrary host file paths. An export returns Markdown to the client,
which chooses where to save it locally.

The node reference is \texttt{pack-id:node-id}. Omitting it uses the
active mission\textquotesingle s next eligible main node, falling back
to its side lane. Manual completion routes remain explicit, even when a
verifier would fail. \texttt{done -\/-complete} logs partial work on
verification failure, returning \texttt{completed: false} and a reason.


\subsubsection{Local lab catalog and submissions}
\texttt{GET /api/lab-catalog} lists the shared open-source references and reviewed
recipes; filters are \texttt{track}, \texttt{level}, \texttt{location} and
\texttt{search}. \texttt{GET /api/lab-catalog/\{id\}} returns a reference or recipe.
Installed pack detail includes \texttt{lab\_references}. These routes perform no
downloads or runtime operations.

\texttt{GET/POST /api/lab-attempts} lists or creates account-scoped records.
Creation accepts \texttt{recipe\_id} and \texttt{offline}; an optional integer
\texttt{case\_key} from 0 to 65535 is allowed only for an offline policy import.
The server assigns a
nonce and the recipe/checker/fixture versions. \texttt{GET /api/lab-attempts/\{id\}}
returns metadata without the nonce. POST actions \texttt{ready}, \texttt{abandon}
and \texttt{submit} require the nonce; submit also requires the assigned versions.
The HTTP body is bounded to 1 MiB, including chunked uploads. The file pilot accepts
an explicit answer schema and evaluates its own rubric. Unknown keys, client pass
claims, cross-account requests and replay are rejected. Local report results have
local authority; neither these records nor server-checked answers issue existing
track proof receipts or XP. The API accepts no local file paths or shell commands.

The policy recipe accepts only \texttt{submission.files.policy.json} as a JSON
object (64 KiB, at most 32 first-match IPv4 rules). Its fixture binds the assigned
case, source, recipe, checker and worker hash. A fresh server-owned worker tests
the declarative policy and returns ``Submission tested (policy model)'', with a
canonical submission hash and counterexample on failure. This is not a live
router or executable-code replay service. Invalid worker output cannot produce a
pass; infrastructure errors leave the attempt retryable. Worker files are
temporary. Manual specialist imports require the exact named observations and
environment identifiers; those remain learner-reported, not server-verified.

Prepared learner-local workspaces submit \texttt{report.evidence} with schema
\texttt{local-artifacts-v1}, recipe/source identity, exact environment fields
and named artifacts containing canonical SHA-256 hashes and JSON observations.
Each observation is bounded to 64 KiB; the complete request remains bounded to
1 MiB. The server derives the required checks and negative controls again.
Missing observations remain unsupported; inconsistent hashes, unexpected fields
and declared pass replacements are rejected. These receipts say ``Local evidence
reviewed'' and retain local authority. They do not authenticate the earlier
environment. Preparation and capture execute only in the learner's companion CLI.

\subsubsection{Playground API}
\begin{verbatim}
GET /api/playground
GET /api/playground/recommendations?minutes=10
GET /api/playground/history
GET /api/playground/insights
POST /api/playground/selection
GET /api/playground/activities/{id}
POST /api/playground/activities/{id}/start
POST /api/playground/activities/{id}/bookmark?enabled=true
GET /api/playground/runs/{run}
POST /api/playground/runs/{run}/{action}
GET /api/nodes/{reference}/checkpoints
POST /api/nodes/{reference}/checkpoints/{checkpoint}
\end{verbatim}
Catalog filters are domain, kind, minutes, topic, difficulty, environment, progress,
format, search and bookmarked. Recommendations also accept comma-separated
exclude IDs. Empty selections include named alternatives. Start accepts
\texttt{deliberate} and \texttt{end\_scored} booleans. Run actions are resume,
save, submit, log, experiment, copy, repair, hint, end, reveal, tour, feedback and
returned. Save, submit and
experiment use a JSON object in \texttt{data}; save optionally sets \texttt{pause}.
Pause applies to examples and free workspaces. Assessed challenges ignore pause
and retain continuous time from their original start until completion or explicit
end. History returns individual runs; the browser groups them by activity ID
without discarding attempts from earlier versions.
Log/example submission use learned and optional next\_step; log also accepts minutes.
Copy accepts reset, blank and fault; repair accepts rule. Resume/copy can explicitly
end a scored attempt with end\_scored. Reveals require a solved or ended attempt.

Tour accepts a zero-based step and status active/finished/skipped, for examples
only. Feedback takes working\_minutes, setup\_minutes, difficulty, clarity,
repetition, environment\_issue and note in data. Difficulty is too-easy/about-right/
too-hard; clarity is clear/unclear; repetition is fresh/welcome-review/unwanted-repeat.
Selection observation takes filters in data and deduplicates by user/day/filter;
GET recommendations does not write observations. Returned records an explicit
return-to-track action. None of these actions awards learning credit. Checkpoint
POST accepts learned and optional minutes and saves a partial entry idempotently.

Every route scopes reads and writes to the authenticated user. Runs never call
track marking. The command endpoint also accepts \texttt{practice} and
\texttt{playground}; remote/local behavior agrees. Stats includes XP by source,
completed practice, solved standalone challenges and saved work. Overview always
retains the next main track task and adds independent practice recommendations.

\subsection{Separate server lab runner}
\path|GET /api/hosted-labs/capabilities| returns configured status, enabled
recipes and runner limits. \path|GET /api/hosted-labs/runs| returns caller-scoped
runs and the latest asynchronous job. \path|POST /api/hosted-labs/runs| accepts
\texttt{recipe\_id} and queues a reviewed recipe launch. The browser polls run
status rather than holding a provisioning request open.
\path|POST /api/hosted-labs/runs/{id}/{action}| supports check, stop, resume,
finish, destroy, bounded save, fixed-target HTTP and restricted router commands.
\path|GET /api/hosted-labs/runs/{id}/files?name=...| reads only selected input
or retained work files. POST wire bodies are limited to 1 MiB; text files and
responses to 256 KiB. A disconnected runner returns 503, leaving local workflows
available. Account scope is selected by the authenticated web dependency; the
browser cannot choose another scope. The provider runtime runs only in the
separate administrative service.
